\documentclass[11pt]{article}

% ─────────────────────────────────────────────────────────────────
% Oracle detector & image-based servoing — projection math
% Self-contained; compile with pdflatex.
% ─────────────────────────────────────────────────────────────────
\usepackage[margin=1in]{geometry}
\usepackage{amsmath}
\usepackage{amssymb}
\usepackage{bm}
\usepackage{booktabs}
\usepackage{xcolor}
\usepackage[colorlinks=true,linkcolor=blue,urlcolor=blue]{hyperref}

\newcommand{\Rcw}{R_{cw}}
\newcommand{\tcw}{\bm{t}_{cw}}
\newcommand{\pd}{\bm{p}_d}
\newcommand{\pt}{\bm{p}_t}
\newcommand{\Pc}{\bm{P}_c}

\title{Oracle Detector and Image-Based Servoing:\\ The Projection Math}
\author{HIL Controller Benchmark}
\date{\today}

\begin{document}
\maketitle

\noindent
The oracle detector and a YOLO-style detector produce the \emph{same kind of
output} (a pixel bounding box). The oracle \emph{computes} that box from
ground-truth positions by running the camera projection \textbf{forward}; the
controller recovers bearing and range by running it \textbf{backward}. This note
derives both directions.

\section{The forward model: the oracle}

The oracle knows two 3D points in the world frame, the drone $\pd$ and the
target $\pt$, plus the drone yaw $\psi$ and pitch $\theta$. It computes where the
target lands on the image, in pixels.

\paragraph{Givens.}
\[
\pd = \begin{bmatrix} x \\ y \\ z \end{bmatrix}_{\text{drone}},\qquad
\pt = \begin{bmatrix} t_x \\ t_y \\ t_z \end{bmatrix}_{\text{target}},\qquad
\text{yaw } \psi,\ \text{pitch } \theta.
\]
Camera intrinsics: $f_x=f_y=554$, principal point $(c_x,c_y)=(320,240)$,
known target height $H=1.5\,\mathrm{m}$.

\subsection*{Step 1 — Translation (world frame)}
Only the position \emph{relative} to the drone matters:
\[
\Delta = \pt - \pd = \begin{bmatrix} dx \\ dy \\ dz \end{bmatrix}.
\]

\subsection*{Step 2 — Rotate world $\to$ body (undo the yaw)}
The body frame is $x$=forward, $y$=left, $z$=up, yawed by $\psi$ about world-$Z$.
Expressing $\Delta$ in body coordinates applies the inverse yaw $R_z(\psi)^\top$:
\[
\begin{bmatrix} \text{fwd} \\ \text{left} \\ \text{up} \end{bmatrix}
=
\underbrace{\begin{bmatrix}
\cos\psi & \sin\psi & 0 \\
-\sin\psi & \cos\psi & 0 \\
0 & 0 & 1
\end{bmatrix}}_{R_z(\psi)^\top}
\begin{bmatrix} dx \\ dy \\ dz \end{bmatrix}.
\]

\subsection*{Step 3 — Apply the pitch (forward--up plane)}
Pitch $\theta$ is a rotation about the body lateral (left) axis:
\[
\begin{bmatrix} \text{fwd}' \\ \text{left} \\ \text{up}' \end{bmatrix}
=
\underbrace{\begin{bmatrix}
\cos\theta & 0 & \sin\theta \\
0 & 1 & 0 \\
-\sin\theta & 0 & \cos\theta
\end{bmatrix}}_{R_{\text{pitch}}(\theta)}
\begin{bmatrix} \text{fwd} \\ \text{left} \\ \text{up} \end{bmatrix}.
\]

\subsection*{Step 4 — Relabel body $\to$ camera-optical axes}
Computer-vision cameras use $Z_c$=forward, $X_c$=right, $Y_c$=down. Converting
from forward/left/up is a signed permutation $P$:
\[
\begin{bmatrix} X_c \\ Y_c \\ Z_c \end{bmatrix}
=
\underbrace{\begin{bmatrix}
0 & -1 & 0 \\
0 & 0 & -1 \\
1 & 0 & 0
\end{bmatrix}}_{P}
\begin{bmatrix} \text{fwd}' \\ \text{left} \\ \text{up}' \end{bmatrix}
=
\begin{bmatrix} -\text{left} \\ -\text{up}' \\ \text{fwd}' \end{bmatrix}.
\]
The visibility guard $Z_c \le 0.1$ means ``behind the camera, not visible.''

\subsection*{The rotation chain collapses to one matrix}
Steps 2--4 fuse into a single world$\to$camera rotation:
\[
\boxed{\;\Rcw = P\,R_{\text{pitch}}(\theta)\,R_z(\psi)^\top,\qquad
\Pc = \begin{bmatrix} X_c \\ Y_c \\ Z_c \end{bmatrix} = \Rcw\,(\pt - \pd)\;}
\]
or, in standard extrinsic $[\,R \mid t\,]$ form,
\[
\Pc = \Rcw\,\pt + \tcw,\qquad \tcw = -\Rcw\,\pd .
\]

\subsection*{Step 5 — Perspective projection (pinhole)}
Flatten the 3D camera point to a pixel with the intrinsic matrix $K$:
\[
s \begin{bmatrix} u \\ v \\ 1 \end{bmatrix}
=
\underbrace{\begin{bmatrix}
f_x & 0 & c_x \\
0 & f_y & c_y \\
0 & 0 & 1
\end{bmatrix}}_{K}
\begin{bmatrix} X_c \\ Y_c \\ Z_c \end{bmatrix},
\qquad s = Z_c .
\]
Dividing by the scale $s=Z_c$ (this division \emph{is} perspective: far things
shrink):
\[
\boxed{\;u = f_x\,\frac{X_c}{Z_c} + c_x,\qquad
v = f_y\,\frac{Y_c}{Z_c} + c_y\;}
\]
This $(u,v)$ is the bounding-box \textbf{center}.

\subsection*{Step 6 — Box size from known height}
An object of physical height $H$ at depth $Z_c$ projects to a pixel height under
the same $1/Z_c$ scaling:
\[
\boxed{\;b_h = \frac{f_y\,H}{Z_c}\;},\qquad b_w = b_h \cdot \text{aspect}.
\]
Box height is \textbf{inversely proportional to depth}: it grows as the drone
nears the target.

\subsection*{The oracle in one line}
\[
s\begin{bmatrix} u \\ v \\ 1 \end{bmatrix}
= K\,P\,R_{\text{pitch}}(\theta)\,R_z(\psi)^\top\,(\pt - \pd),
\qquad b_h = \frac{f_y H}{Z_c}.
\]

\section{The inverse model: what a detector gives the controller}

A single detection contains a class label, a confidence, a pixel center
$(c_x,c_y)$ and a pixel size $(b_w,b_h)$ --- \emph{not} a 3D world position. That
pixel data splits into two physical meanings, each obtained by inverting a step
above.

\paragraph{(a) Box center $\to$ bearing.}
Inverting Step 5 turns a pixel offset into an angle (direction to steer):
\[
\frac{X_c}{Z_c} = \frac{u - c_x}{f_x} \approx \tan(\text{horizontal angle}),
\qquad
\frac{Y_c}{Z_c} = \frac{v - c_y}{f_y}.
\]
The controller normalises this to a centring error in $[-1,1]$ (zero when
centred):
\[
e_x = \frac{c_x - W/2}{W/2},\qquad e_y = \frac{c_y - H_{\text{img}}/2}{H_{\text{img}}/2}.
\]

\paragraph{(b) Box size $\to$ range.}
Inverting Step 6 (the oracle's box equation solved for depth) recovers $Z_c$ from
the known height $H$:
\[
\boxed{\;Z_c = \frac{f_y\,H}{b_h}\;}.
\]
A bigger box means a smaller $Z_c$: ``you are close.'' A proportional law may use
the dimensionless ratio $b_h/H_{\text{img}}$ as a range proxy instead.

\section{Q\&A: Clarifications}

\subsection*{Q: Yaw of who and pitch of who?}

Both are the \textbf{drone's} yaw and pitch --- not the target's.

\begin{itemize}
    \item \textbf{Yaw} $\psi$: which compass direction the drone's nose is pointing.
    \item \textbf{Pitch} $\theta$: how much the drone's nose is tilted up or down.
\end{itemize}

Steps 2 and 3 undo the drone's own orientation so you can describe where the
target sits relative to the camera. The target's yaw (which way the stop sign
faces) is available in \texttt{/sim/target\_pose} but is deliberately discarded:

\begin{verbatim}
x, y, z, pitch, yaw = self.drone   # drone orientation used
tx, ty, tz, _       = self.target  # target yaw thrown away
\end{verbatim}

It does not matter which way the sign is facing --- only where it is in space.

\subsection*{Q: What are camera intrinsics and how are they used?}

Camera intrinsics are \textbf{fixed properties of the lens and sensor} --- they
never change during flight. They answer one question: for this specific camera,
how do distances in the 3D world map to distances in pixels?

There are four numbers:

\begin{center}
\begin{tabular}{@{}clll@{}}
\toprule
Parameter & Value (sim) & Meaning \\
\midrule
$f_x$ & 554\,px & horizontal ``zoom'' \\
$f_y$ & 554\,px & vertical ``zoom'' \\
$c_x$ & 320\,px & pixel column of image centre \\
$c_y$ & 240\,px & pixel row of image centre \\
\bottomrule
\end{tabular}
\end{center}

\paragraph{What focal length actually means.}
$f_x = 554$ pixels is not 554\,mm like a camera lens. It is the focal length
expressed in pixel units, which encodes the field of view:
\[
f_x = \frac{W}{2\tan(\mathrm{HFOV}/2)}.
\]
The sim has $\mathrm{HFOV}=60°$ and $W=640$, so:
\[
f_x = \frac{640}{2\tan(30°)} = \frac{640}{1.155} \approx 554.
\]
A larger $f_x$ means a narrower field of view (more zoom); a smaller $f_x$ means
a wider field of view (fisheye-like).

\paragraph{How they are used in Step 5.}
After all the rotations you have $(X_c, Y_c, Z_c)$ in metres. The intrinsics
convert that to pixels:
\[
u = f_x \cdot \frac{X_c}{Z_c} + c_x.
\]
Breaking this apart:
\begin{itemize}
    \item $X_c / Z_c$ is a \textbf{pure angle} (metres/metres, dimensionless).
          It is the same regardless of which camera you have.
    \item Multiply by $f_x$ to \textbf{scale that angle into pixels} for this
          specific camera's sensor.
    \item Add $c_x = 320$ to \textbf{shift the origin} from image centre to the
          top-left corner, because pixels are counted from the top-left.
\end{itemize}

\paragraph{Numerical example.}
Target 1\,m to the right at 10\,m depth:
\[
u = 554 \cdot \frac{1}{10} + 320 = 55.4 + 320 = 375\,\text{px} \quad
\text{(55 pixels right of centre --- makes sense)}.
\]

\paragraph{The principal point $(c_x, c_y)$.}
This is where a perfectly centred target ($X_c = Y_c = 0$, dead ahead) lands on
the image. Ideally it is the exact middle:
\[
c_x = W/2 = 320, \qquad c_y = H_\text{img}/2 = 240.
\]
On a real physical camera it is slightly off-centre due to manufacturing
tolerances and is measured during camera calibration. In the sim it is exactly
centred.

\paragraph{Intrinsics vs.\ extrinsics.}

\begin{center}
\begin{tabular}{@{}lll@{}}
\toprule
 & What it answers & Changes during flight? \\
\midrule
\textbf{Extrinsics} (Steps 2--4) & Where is the target relative to my camera \emph{right now}? & Yes, every frame \\
\textbf{Intrinsics} (Step 5)     & Given its position, what pixel does it land on?              & Never \\
\bottomrule
\end{tabular}
\end{center}

\noindent
The intrinsics are loaded once at startup and shared identically by the oracle,
the FSM, and the IBVS controller --- because all three must agree on what a
given pixel coordinate means in the real world.

\section{The punchline}

The oracle and the detector are the same projection equation run in opposite
directions:

\begin{center}
\begin{tabular}{@{}llll@{}}
\toprule
 & Input & Output & Direction \\
\midrule
Oracle      & known 3D positions & pixel box $(u,v,b_h)$ & world $\to$ image (forward) \\
Detector    & the image          & pixel box $(u,v,b_h)$ & pixels $\to$ box (measured) \\
Controller  & pixel box          & bearing $e_x,e_y$ + depth $Z_c$ & image $\to$ world (inverse) \\
\bottomrule
\end{tabular}
\end{center}

\noindent
Because the oracle emits the \emph{same} pixel box a real detector would, it is a
drop-in replacement on the detection topic. The box \textbf{center is a
direction} and the box \textbf{size is a distance}; together they form a ray plus
a range --- a full relative position, if one ever wants to reconstruct it. Basic
image-based servoing never does: it simply drives ``centre the box and make it
the right size.''

\end{document}